\section{Precisão dos contrastes de frequência com ruído fixado}
\label{sec:contrastes-condicionais}

A comparação com marginalização dos parâmetros de ruído, apresentada
no Capítulo~\ref{chap:compressao}, deixou intervalos numéricos amplos
para as diferenças de limites superiores. Para examinar a aproximação
de frequência em uma situação mais simples, consideram-se aqui as
primeiras 50 realizações da priori da geometria de 16 pulsares e oito
frequências. A escolha dos identificadores precedeu o cálculo dos
contrastes. Os parâmetros de amplitude e ruído são fixados nos valores
de geração, e somente $u$ é integrado. A pergunta passa a ser quanto
a resposta de frequência altera o quantil de massa quando esses
parâmetros são conhecidos.

As duas representações positivas da densidade posterior, construídas
com as quadraturas de ordens 32 e 64, fornecem CDFs $F_{32}$ e $F_{64}$.
A margem de sensibilidade vertical $\varepsilon_F=0{,}002$, usada neste
capítulo, é convertida em um intervalo horizontal para o quantil:
\begin{equation}
 I_{95}=\left[
 \min_{j\in\{32,64\}}F_j^{-1}(0{,}95-\varepsilon_F),\;
 \max_{j\in\{32,64\}}F_j^{-1}(0{,}95+\varepsilon_F)
 \right].
\end{equation}
Essa inversão incorpora a inclinação local da CDF: uma mesma incerteza
vertical pode corresponder a larguras distintas em massa. Para um
contraste $Q^{(X)}_{95}-Q^{(Y)}_{95}$, o intervalo é
$[I^{(X)}_- - I^{(Y)}_+,\ I^{(X)}_+ - I^{(Y)}_-]$.
A integração independente de cada segmento polinomial por
Gauss--Legendre confirmou as 900 inversões utilizadas, com resíduo
máximo de $8{,}7\times10^{-15}$ na CDF. A incerteza relevante continua
sendo a representação da posterior, e não a solução da equação inversa.

\begin{table}[htbp]
\centering
\caption{Contrastes condicionais de $Q_{95}(u)$ nas 50 realizações.
O intervalo da média incorpora a sensibilidade numérica de todos os
pares; $r_{\max}$ é o maior afastamento de um extremo em relação ao
contraste nominal individual.}
\label{tab:contrastes-condicionais}
\small
\begin{tabular}{lrrr}
\toprule
Contraste & Média & Intervalo numérico da média & $r_{\max}$\\
\midrule
$C_{\beta,\mathrm G}-B_{\mathrm G}$ & $-0{,}000215$ & $[-0{,}004190;\ 0{,}003760]$ & $0{,}006940$\\
$C_{\mathrm{full},\mathrm G}-C_{\beta,\mathrm G}$ & $0{,}00000007$ & $[-0{,}003982;\ 0{,}003983]$ & $0{,}007122$\\
$C_{\beta,\mathrm{CN}}-B_{\mathrm{CN}}$ & $-0{,}000299$ & $[-0{,}004351;\ 0{,}003754]$ & $0{,}006927$\\
$C_{\mathrm{full},\mathrm{CN}}-C_{\beta,\mathrm{CN}}$ & $-0{,}00000029$ & $[-0{,}004064;\ 0{,}004063]$ & $0{,}007105$\\
\bottomrule
\end{tabular}
\end{table}

A sensibilidade numérica de cada contraste é inferior a $0{,}01$ em
cada direção. Para $C_\beta-B$, 49 dos 50 intervalos individuais ficam
inteiramente em $[-0{,}01;0{,}01]$, tanto no controle gaussiano quanto
nas estatísticas quadráticas físicas. Para $C_{\mathrm{full}}-C_\beta$,
essa condição vale nos 50 pares de cada família. Os intervalos das
quatro médias também cabem nessa faixa. Portanto, a aproximação produz
pequenos deslocamentos médios do limite superior nessa amostra
condicional, embora um dos pares de cada contraste $C_\beta-B$ ainda
ultrapasse a janela de equivalência.

Os intervalos da tabela descrevem a sensibilidade às representações
numéricas e à margem vertical especificada. A dispersão entre
realizações constitui outra fonte de incerteza; não foi usada aqui
para uma afirmação de equivalência populacional. Além disso, os
contrastes CN permanecem condicionais à aproximação normal da
verossimilhança. A conclusão se refere à inferência de massa com ruído
fixado nessa geometria, enquanto o efeito conjunto da aproximação e
da marginalização permanece a questão aberta do estudo anterior.
